% GENERATED FILE. Edit canonical lessons, then run npm run generate:books.
% canonical-source-hash: fnv1a64:010dac6be3237f6f
% canonical-lessons: JA-C220-kenkou, JA-C220-joubu, JA-C220-futoru, JA-C220-yaseru, JA-C220-shokuyoku

\chapter{Weight and Staying Healthy}
\label{ch:ja-weight-and-staying-healthy}
\hlchaptermodality{220}

\begin{chapteropening}
\textbf{By the end of this chapter:} \emph{I can talk about health, weight and appetite.}

Complete the last lesson of chapter 220: I can talk about health, weight and appetite.
\end{chapteropening}

\section[\texorpdfstring{\ja{けんこう} (kenkō)}{けんこう (kenkō)}]{\ja{けんこう} (kenkō) — health}

\label{lesson:JA-C220-kenkou}



\begin{quote}
Before the new one: say the Japanese for to collapse, then the Japanese for an injury.
\end{quote}

\subsection*{You'll want to know: \ja{けんこう}}
\textbf{\ja{けんこう}} — \emph{kenkō} — \textquotedblleft{}health\textquotedblright{}.

\begin{quote}
\textbf{\ja{けんこうは} \ja{たいせつです。}} — \emph{kenkō wa taisetsu desu} — \textquotedblleft{}Health is important.\textquotedblright{}
\end{quote}

\begin{cousinweb}
\textbf{\ja{けんこう}} (\emph{kenkō}) is Sino-Japanese (kango): two Chinese characters read with their on'yomi, readings borrowed from Chinese. They are \emph{ken}, \textquotedblleft{}strong\textquotedblright{} (Middle Chinese \emph{gjonH}) and \emph{kō}, \textquotedblleft{}at ease\textquotedblright{} (\emph{khang}) (KANJIDIC2; Middle Chinese in Baxter's transcription).
\end{cousinweb}

\begin{grammarlens}[title={What you've built}]
One more word for everyday life at A2.
\end{grammarlens}

\subsection*{Your turn}
\begin{itemize}
\raggedright
  \item \emph{Say it:} \emph{kenkō}
  \item \emph{Say it:} \emph{kenkō}, once more
  \item \emph{Say it:} \emph{kenkō wa taisetsu desu} — Health is important.
  \item \emph{Recall:} say the Japanese for to collapse, then the Japanese for an injury, then say \emph{kenkō} again
\end{itemize}

\begin{tcolorbox}[breakable,colback=teal!4,colframe=teal!35!black,title={Before you move on}]
What does \ja{けんこう} mean? (\textquotedblleft{}Health\textquotedblright{}.) Say it once more.
\end{tcolorbox}
\section[\texorpdfstring{\ja{じょうぶ} (jōbu)}{じょうぶ (jōbu)}]{\ja{じょうぶ} (jōbu) — robust; sturdy}

\label{lesson:JA-C220-joubu}



\begin{quote}
Before the new one: say the Japanese for an injury, then the Japanese for health.
\end{quote}

\subsection*{You'll want to know: \ja{じょうぶ}}
\textbf{\ja{じょうぶ}} — \emph{jōbu} — \textquotedblleft{}robust; sturdy\textquotedblright{}.

\begin{quote}
\textbf{\ja{ちちは} \ja{からだが} \ja{じょうぶです。}} — \emph{chichi wa karada ga jōbu desu} — \textquotedblleft{}My father has a strong body.\textquotedblright{}
\end{quote}

\begin{cousinweb}
\textbf{\ja{じょうぶ}} (\emph{jōbu}) is Sino-Japanese (kango): two Chinese characters read with their on'yomi, readings borrowed from Chinese. They are \emph{jō}, \textquotedblleft{}ten feet, a full height\textquotedblright{} (Middle Chinese \emph{drjangX}) and \emph{bu}, \textquotedblleft{}man\textquotedblright{} (\emph{pju}) (KANJIDIC2; Middle Chinese in Baxter's transcription).
\end{cousinweb}

\begin{grammarlens}[title={What you've built}]
One more word for everyday life at A2.
\end{grammarlens}

\subsection*{Your turn}
\begin{itemize}
\raggedright
  \item \emph{Say it:} \emph{jōbu}
  \item \emph{Say it:} \emph{jōbu}, once more
  \item \emph{Say it:} \emph{chichi wa karada ga jōbu desu} — My father has a strong body.
  \item \emph{Recall:} say the Japanese for an injury, then the Japanese for health, then say \emph{jōbu} again
\end{itemize}

\begin{tcolorbox}[breakable,colback=teal!4,colframe=teal!35!black,title={Before you move on}]
What does \ja{じょうぶ} mean? (\textquotedblleft{}Robust; sturdy\textquotedblright{}.) Say it once more.
\end{tcolorbox}
\section[\texorpdfstring{\ja{ふとる} (futoru)}{ふとる (futoru)}]{\ja{ふとる} (futoru) — to gain weight}

\label{lesson:JA-C220-futoru}



\begin{quote}
Before the new one: say the Japanese for health, then the Japanese for robust.
\end{quote}

\subsection*{You'll want to know: \ja{ふとる}}
\textbf{\ja{ふとる}} — \emph{futoru} — \textquotedblleft{}to gain weight\textquotedblright{}.

\begin{quote}
\textbf{\ja{いぬは} \ja{ふとりましたか。}} — \emph{inu wa futorimashita ka} — \textquotedblleft{}Has the dog put on weight?\textquotedblright{}
\end{quote}

Polite \textbf{\ja{ふとります}}, \emph{futorimasu}; a godan verb.

\begin{cousinweb}
\textbf{\ja{ふとる}} (\emph{futoru}) is a native Japanese word (wago), inherited rather than borrowed and built on the same stem as \emph{futoi}, \textquotedblleft{}thick\textquotedblright{}; its kanji was attached later as a kun'yomi (JMdict, KANJIDIC2).
\end{cousinweb}

\begin{grammarlens}[title={What you've built}]
One more word for everyday life at A2.
\end{grammarlens}

\subsection*{Your turn}
\begin{itemize}
\raggedright
  \item \emph{Say it:} \emph{futoru}
  \item \emph{Say it:} \emph{futoru}, once more
  \item \emph{Say it:} \emph{inu wa futorimashita ka} — Has the dog put on weight?
  \item \emph{Recall:} say the Japanese for health, then the Japanese for robust, then say \emph{futoru} again
\end{itemize}

\begin{tcolorbox}[breakable,colback=teal!4,colframe=teal!35!black,title={Before you move on}]
What does \ja{ふとる} mean? (\textquotedblleft{}To gain weight\textquotedblright{}.) Say it once more.
\end{tcolorbox}
\section[\texorpdfstring{\ja{やせる} (yaseru)}{やせる (yaseru)}]{\ja{やせる} (yaseru) — to lose weight; to get thin}

\label{lesson:JA-C220-yaseru}



\begin{quote}
Before the new one: say the Japanese for robust, then the Japanese for to gain weight.
\end{quote}

\subsection*{You'll want to know: \ja{やせる}}
\textbf{\ja{やせる}} — \emph{yaseru} — \textquotedblleft{}to lose weight; to get thin\textquotedblright{}.

\begin{quote}
\textbf{\ja{ちちは} \ja{やせませんでした。}} — \emph{chichi wa yasemasen deshita} — \textquotedblleft{}My father did not lose weight.\textquotedblright{}
\end{quote}

Polite \textbf{\ja{やせます}}, \emph{yasemasu}; an ichidan verb.

\begin{cousinweb}
\textbf{\ja{やせる}} (\emph{yaseru}) is a native Japanese word (wago), inherited rather than borrowed from Chinese; its kanji was attached later as a kun'yomi, and its older origin is not settled (JMdict, KANJIDIC2).
\end{cousinweb}

\begin{grammarlens}[title={What you've built}]
One more word for everyday life at A2.
\end{grammarlens}

\subsection*{Your turn}
\begin{itemize}
\raggedright
  \item \emph{Say it:} \emph{yaseru}
  \item \emph{Say it:} \emph{yaseru}, once more
  \item \emph{Say it:} \emph{chichi wa yasemasen deshita} — My father did not lose weight.
  \item \emph{Recall:} say the Japanese for robust, then the Japanese for to gain weight, then say \emph{yaseru} again
\end{itemize}

\begin{tcolorbox}[breakable,colback=teal!4,colframe=teal!35!black,title={Before you move on}]
What does \ja{やせる} mean? (\textquotedblleft{}To lose weight; to get thin\textquotedblright{}.) Say it once more.
\end{tcolorbox}
\section[\texorpdfstring{\ja{しょくよく} (shokuyoku)}{しょくよく (shokuyoku)}]{\ja{しょくよく} (shokuyoku) — an appetite}

\label{lesson:JA-C220-shokuyoku}



\begin{quote}
Before the new one: say the Japanese for to gain weight, then the Japanese for to lose weight.
\end{quote}

\subsection*{You'll want to know: \ja{しょくよく}}
\textbf{\ja{しょくよく}} — \emph{shokuyoku} — \textquotedblleft{}an appetite\textquotedblright{}.

\begin{quote}
\textbf{\ja{しょくよくは} \ja{ありますか。}} — \emph{shokuyoku wa arimasu ka} — \textquotedblleft{}Do you have an appetite?\textquotedblright{}
\end{quote}

\begin{cousinweb}
\textbf{\ja{しょくよく}} (\emph{shokuyoku}) is Sino-Japanese (kango): two Chinese characters read with their on'yomi, readings borrowed from Chinese. They are \emph{shoku}, \textquotedblleft{}to eat\textquotedblright{} (Middle Chinese \emph{zyik}) and \emph{yoku}, \textquotedblleft{}desire\textquotedblright{} (\emph{yowk}) (KANJIDIC2; Middle Chinese in Baxter's transcription).
\end{cousinweb}

\begin{grammarlens}[title={What you've built}]
One more word for everyday life at A2.
\end{grammarlens}

\subsection*{Your turn}
\begin{itemize}
\raggedright
  \item \emph{Say it:} \emph{shokuyoku}
  \item \emph{Say it:} \emph{shokuyoku}, once more
  \item \emph{Say it:} \emph{shokuyoku wa arimasu ka} — Do you have an appetite?
  \item \emph{Recall:} say the Japanese for to gain weight, then the Japanese for to lose weight, then say \emph{shokuyoku} again
\end{itemize}

\begin{tcolorbox}[breakable,colback=teal!4,colframe=teal!35!black,title={Before you move on}]
What does \ja{しょくよく} mean? (\textquotedblleft{}An appetite\textquotedblright{}.) Say it once more.
\end{tcolorbox}
